\begin{examplebox}

	\textbf{\textcolor{CExample}{例 1（基础）}}

	\textbf{\textcolor{CExample}{题目：}}
	已知双曲线 $\frac{x^2}{16}-\frac{y^2}{9}=1$，点 $P\left(5,\frac{9}{4}\right)$ 在双曲线上，求 $|PF_1|$ 和 $|PF_2|$。

	\textbf{\textcolor{CExample}{解题步骤：}}
	\begin{description}[style=nextline, leftmargin=2.8em, labelsep=0.5em, itemsep=0.45em]
		\item[\textbf{第一步（求参数）：}] 由方程得 $a=4$, $b=3$, $c=5$, $e=\frac{5}{4}$。
			\par\textbf{理由：} 在标准方程 $\frac{x^2}{a^2}-\frac{y^2}{b^2}=1$ 中，$a^2=16$, $b^2=9$，$c=\sqrt{a^2+b^2}=5$。
		\item[\textbf{第二步（判断支）：}] 因为 $x_0=5\ge a=4$，所以点 $P$ 在右支。
			\par\textbf{理由：} 双曲线上点横坐标满足 $|x_0|\ge a$，右支对应 $x_0\ge a$。
		\item[\textbf{第三步（套公式）：}] 在右支，焦半径公式为 $|PF_1|=e x_0 + a$，$|PF_2|=e x_0 - a$。代入数值：
			\[
				\begin{aligned}
					|PF_1| &= e x_0 + a \\
					&= \frac{5}{4}\times5+4 \\
					&= \frac{41}{4},\\
					|PF_2| &= e x_0 - a \\
					&= \frac{5}{4}\times5-4 \\
					&= \frac{9}{4}.
			\end{aligned}
			\]
			\par\textbf{理由：} 右支时焦半径公式去绝对值，代入 $e=\frac{5}{4}$, $x_0=5$, $a=4$ 计算。
	\end{description}

	\textbf{\textcolor{CExample}{关键结论：}}
	$|PF_1|=\frac{41}{4}$, $|PF_2|=\frac{9}{4}$。

	\medskip
	\textbf{\textcolor{CExample}{例 2（稍变形）}}

	\textbf{\textcolor{CExample}{题目：}}
	焦点在 $x$ 轴的双曲线右支上一点 $P(5,y_0)$ 到两焦点的距离分别为 $13$ 和 $3$，求该双曲线的标准方程。

	\textbf{\textcolor{CExample}{解题步骤：}}
	\begin{description}[style=nextline, leftmargin=2.8em, labelsep=0.5em, itemsep=0.45em]
		\item[\textbf{第一步（设标准方程）：}] 设双曲线方程为 $\frac{x^2}{a^2}-\frac{y^2}{b^2}=1$（$a>0,b>0$），离心率 $e=\frac{c}{a}$。
			\par\textbf{理由：} 焦点在 $x$ 轴的标准形式。
		\item[\textbf{第二步（列焦半径方程组）：}] 因为 $P$ 在右支，根据焦半径公式有
			\[
				\begin{cases}
					e\cdot5 + a = 13,\\
					e\cdot5 - a = 3.
			\end{cases}
			\]
			\par\textbf{理由：} 右支时 $|PF_1|=e x_0 + a$，$|PF_2|=e x_0 - a$。
		\item[\textbf{第三步（解出 $a$ 和 $e$）：}] 两式相加得 $10e=16$，进而：
			\[
				\begin{aligned}
					10e &= 16,\\
					e &= \frac{8}{5},\\
					a &= 13 - 5e \\
					&= 13 - 8 \\
					&= 5.
			\end{aligned}
			\]
			\par\textbf{理由：} 加减消元，代入运算。
		\item[\textbf{第四步（求 $b$）：}]
			\[
				\begin{aligned}
					c &= ae \\
					&= 5 \times \frac{8}{5} \\
					&= 8,\\
					b^2 &= c^2 - a^2 \\
					&= 64 - 25 \\
					&= 39.
			\end{aligned}
			\]
			\par\textbf{理由：} 双曲线参数关系 $c^2=a^2+b^2$。
	\end{description}

	\textbf{\textcolor{CExample}{关键结论：}}
	双曲线方程为 $\frac{x^2}{25}-\frac{y^2}{39}=1$。

	\medskip
	\textbf{\textcolor{CExample}{例 3（综合应用）}}

	\textbf{\textcolor{CExample}{题目：}}
	设 $F_1,F_2$ 是双曲线 $\frac{x^2}{16}-\frac{y^2}{9}=1$ 的左、右焦点，点 $P$ 在双曲线上，且 $|PF_1|:|PF_2|=3:1$，求点 $P$ 的坐标。

	\textbf{\textcolor{CExample}{解题步骤：}}
	\begin{description}[style=nextline, leftmargin=2.8em, labelsep=0.5em, itemsep=0.45em]
		\item[\textbf{第一步（求参数）：}] 由双曲线方程得 $a=4$, $c=5$, $e=\frac{5}{4}$。
			\par\textbf{理由：} 标准方程系数。
		\item[\textbf{第二步（设点于右支列方程）：}] 由于 $|PF_1|>|PF_2|$，故点 $P$ 应在右支。右支时 $|PF_1|=e x_0+a$, $|PF_2|=e x_0-a$，比例条件给出
			\[
				e x_0+a = 3(e x_0-a).
			\]
			\par\textbf{理由：} 焦半径公式和比例关系。
		\item[\textbf{第三步（解横坐标）：}] 化简方程得
			\[
				\begin{aligned}
					e x_0 + a &= 3e x_0 - 3a \\
					4a &= 2e x_0 \\
					x_0 &= \frac{2a}{e} \\
					&= \frac{2a}{c/a} \\
					&= \frac{2a^2}{c} \\
					&= \frac{2\times16}{5} \\
					&= \frac{32}{5} \\
					&= 6.4.
			\end{aligned}
			\]
			\par\textbf{理由：} 代数运算，注意 $e=c/a$。
		\item[\textbf{第四步（求纵坐标）：}] 将 $x_0$ 代入双曲线方程：
			\[
				\begin{aligned}
					\frac{\left(\frac{32}{5}\right)^2}{16} - \frac{y_0^2}{9} &= 1 \\
					\frac{64}{25} - \frac{y_0^2}{9} &= 1 \\
					\frac{y_0^2}{9} &= \frac{64}{25} - 1 \\
					\frac{y_0^2}{9} &= \frac{39}{25} \\
					y_0^2 &= \frac{351}{25} \\
					y_0 &= \pm \frac{3\sqrt{39}}{5}.
			\end{aligned}
			\]
			\par\textbf{理由：} 代入化简，解出 $y_0^2$ 并开方。
	\end{description}

	\textbf{\textcolor{CExample}{关键结论：}}
	点 $P$ 的坐标为 $\left(\frac{32}{5},\pm\frac{3\sqrt{39}}{5}\right)$（右支）。左支无解（$x_0=6.4$ 不满足 $x_0\le -4$）。

\end{examplebox}
